\chapter{Console Commands}\label{app:console}

The console is on the USB port (type \cmd{menu} on a line by itself) and on
TCP port 9300. Each command ends with a line that starts with \code{ok} or
\code{error:}. When the console opens, it shows the firmware version, the
model name, the IP address, the active links, and the source of the
settings.

\begin{longtable}{@{}P{0.27\linewidth}P{0.68\linewidth}@{}}
\caption{Console commands}\label{tab:console}\\
\toprule
\textbf{Command} & \textbf{Function} \\
\midrule
\endfirsthead
\toprule
\textbf{Command} & \textbf{Function} \\
\midrule
\endhead
\bottomrule
\endfoot
\cmd{help}, \cmd{menu} & Show the list of commands. \\
\cmd{status} & Show the state, the position in mm, the pen, the queue depth,
the counters, and the data of the links (table~\ref{tab:status}). \\
\cmd{config} & Show all settings with their values, their descriptions, and
when they take effect. \\
\cmd{set} \textit{key value} & Make sure that the value is correct, apply it,
and save it in the nonvolatile memory. Appendix~\ref{app:settings} lists the
keys. \\
\cmd{defaults} & Set all settings to the factory values. The motors must be
disabled. \\
\cmd{enable} & Enable the motors. This also clears a latched fault. With
\key{hlfb_check}~1, the board waits for HLFB from all servos. \\
\cmd{disable} & Disable the motors. A job that is in the queue is dropped. \\
\cmd{abort} & Stop now, drop the job, and raise the pen. \\
\cmd{pause} & Stop at the end of the segments in the queue, and raise the pen. \\
\cmd{resume} & Lower the pen and continue after \cmd{pause}. \\
\cmd{origin} & The present position becomes plotter 0,0. No job can be
plotting. \\
\cmd{goto} \textit{x\_mm y\_mm} & Move to the position with the pen up. No job
can be plotting. The position stays in the plot area. \\
\cmd{pen up}, \cmd{pen down} & Raise or lower the pen to test it. No job can
be plotting. \\
\cmd{clear} & Clear a latched fault. Remove the cause first. \\
\cmd{diag} & Show the motion diagnostics, the last job of the optimizer, and
the data of the USB link. \\
\cmd{diag reset} & Set the timing counters to zero, then show the diagnostics. \\
\cmd{tel} \textit{ms} & TCP console only. Send a line \code{TEL ms= x= y= pen=
state= queue=} every \textit{ms} milliseconds (10 or more; 0 stops it). A line
is sent only when the TCP buffer has space, so a slow client cannot stop the
motion. The desktop app uses 25~ms. \\
\cmd{restart} & Restart the board. \\
\cmd{version} & Show the name and the version of the firmware. \\
\cmd{exit}, \cmd{quit}, \cmd{q} & Close the console. On USB, the port goes back
to HP-GL. \\
\end{longtable}

\section*{Error Messages}

\begin{table}[htbp]
\centering
\caption{Common console errors}\label{tab:console-errors}
\begin{tabularx}{\linewidth}{@{}P{0.45\linewidth}L@{}}
\toprule
\textbf{Message} & \textbf{Meaning} \\
\midrule
\code{unknown key (try: config)} & The setting name is not correct. \\
\code{disable motors first} & The setting needs the motors disabled. \\
\code{a job is running; change it when the plotter is idle} & The setting can change only between jobs. \\
\code{applied but NVM write failed} & The board uses the value but could not save it. The supply voltage is too low. \\
\code{plotter busy} & The command needs the machine to be idle. \\
\code{unknown command (try: help)} & The command name is not correct. \\
\bottomrule
\end{tabularx}
\end{table}
